% demo-deck.tex
% Synthetic demo deck for latex-defense-zh. The deck covers the twelve layouts
% in references/slide-layouts.md. A "% layout: <name>" line precedes each frame.
% All names, titles, numbers, and references in this file are synthetic.
%
% Default theme (YanshanDefense):
%   latexmk -xelatex demo-deck.tex
% Generic theme:
%   latexmk -xelatex -jobname=demo-generic "-usepretex=\def\DefenseThemeName{GenericDefense}" demo-deck.tex
\providecommand{\DefenseThemeName}{YanshanDefense}
\documentclass[aspectratio=169,11pt]{ctexbeamer}
\usetheme{\DefenseThemeName}

\DefenseSetup{logo=logo.png,stage=predefense}
\title{面向合成数据的\\示例建模方法研究}
\author{示例作者}
\DefenseSupervisor{示例导师\quad 教授}
\DefenseSubject{计算机科学与技术}
\DefenseSchool{示例学院}
\date{2026年9月}

\DefenseAddChapter{1}{绪论}
\DefenseAddChapter{2}{理论基础与问题框架}
\DefenseAddChapter{3}{示例方法一研究}
\DefenseAddChapter{4}{示例方法二研究}
\DefenseAddChapter{5}{示例平台设计与应用}
\DefenseAddChapter{6}{总结与展望}

\DefenseDefineLabel{fig:demo-model}{3-1}
\DefenseDefineLabel{eq:demo-output}{3-1}
\DefenseDefineLabel{eq:demo-loss}{3-2}

\begin{document}

% layout: cover
\begin{frame}[plain]
\DefenseCoverFrame
\end{frame}

% layout: toc
\begin{frame}[plain]
\DefenseTocFrame{0}
\end{frame}

% layout: bullets
\begin{frame}{1.1 研究背景与意义}
\DefenseSubsection{1.1.1 示例任务背景}
\begin{itemize}
  \item 示例任务的数据规模持续增长，人工分析的成本随之上升。
  \item 现有流程依赖人工经验，分析结果的\DefenseHighlight{一致性}不足。
  \item 示例任务需要一种可复现、可更新的自动建模方法。
\end{itemize}
\DefenseTakeaway{示例任务需要可复现、可更新的自动建模方法}
\end{frame}

% layout: figure-bullets
\begin{frame}{1.2 国内外研究现状}
\DefenseSubsection{1.2.1 数据驱动建模方法}
\begin{columns}[T,onlytextwidth]
  \column{0.54\textwidth}
  \DefenseFigure[height=0.5\textheight]{example-image-a}{图1-1\quad 示例方法类别统计}
  \column{0.42\textwidth}
  \begin{itemize}
    \item 规则方法可解释性强，但覆盖范围有限。
    \item 统计学习方法依赖人工特征设计。
    \item 深度学习方法精度较高，但需要大量标注样本。
  \end{itemize}
\end{columns}
\DefenseTakeaway{现有方法难以同时兼顾精度与标注成本}
\end{frame}

% layout: cards
\begin{frame}{1.3 现存问题与挑战}
\DefenseCard{问题一}{示例数据存在\DefenseHighlight{分布差异}，单一模型难以覆盖全部场景。}
\DefenseCard{问题二}{标注样本\DefenseHighlight{稀缺}，监督学习方法的泛化能力不足。}
\DefenseCard{问题三}{模型部署后缺少\DefenseHighlight{在线更新}机制，精度随时间下降。}
\end{frame}

% layout: figure
\begin{frame}{1.4 研究内容与章节安排}
\DefenseFigure{example-image-16x9}{图1-2\quad 示例论文组织结构}
\DefenseTakeaway{第 3、4 章分别回答问题一与问题二，第 5 章完成平台应用}
\end{frame}

% layout: toc
\begin{frame}[plain]
\DefenseTocFrame{3}
\end{frame}

% layout: figure-bullets
\begin{frame}{3.1 引言}
\DefenseSubsection{研究内容 1：示例方法一研究}
\DefenseFigure[height=0.36\textheight]{example-image-16x9}{图3-1\quad 示例技术路线}
\begin{itemize}
  \item 首先构建示例模型，其次设计训练目标，最后在合成数据上验证。
\end{itemize}
\DefenseTakeaway{本章针对问题一提出示例方法一}
\end{frame}

% layout: equations-figure
\begin{frame}{3.3 示例模型建立}
\DefenseSubsection{3.3.1 模型结构与训练目标}
\begin{columns}[T,onlytextwidth]
  \column{0.52\textwidth}
  \begin{DefenseSource}
  \begin{align*}
    \hat{y}_t &= f_{\theta}(\mathbf{x}_t) + b \tag*{(3-1)}\\
    \mathcal{L}(\theta) &= \frac{1}{N}\sum_{t=1}^{N}\left(y_t-\hat{y}_t\right)^2 \tag*{(3-2)}
  \end{align*}
  式~\eqref{eq:demo-loss} 为训练目标，模型结构见图~\ref{fig:demo-model}\cite{demo-ref}。
  \end{DefenseSource}
  \begin{itemize}
    \item 以均方误差作为训练目标。
  \end{itemize}
  \column{0.44\textwidth}
  \DefenseFigure[height=0.45\textheight]{example-image-b}{图3-1\quad 示例模型结构}
\end{columns}
\DefenseTakeaway{示例模型以均方误差为训练目标}
\end{frame}

% layout: figure-grid
\begin{frame}{3.4 案例研究与结果分析}
\DefenseSubsection{3.4.1 示例实验一}
\begingroup\centering
\begin{minipage}[t]{0.32\linewidth}
  \DefenseFigure[height=0.21\textheight]{example-image-a}{(a) 方法 A 结果}
\end{minipage}\hspace{0.08\linewidth}%
\begin{minipage}[t]{0.32\linewidth}
  \DefenseFigure[height=0.21\textheight]{example-image-b}{(b) 方法 B 结果}
\end{minipage}\par\vspace{2pt}
\begin{minipage}[t]{0.32\linewidth}
  \DefenseFigure[height=0.21\textheight]{example-image-c}{(c) 方法 C 结果}
\end{minipage}\hspace{0.08\linewidth}%
\begin{minipage}[t]{0.32\linewidth}
  \DefenseFigure[height=0.21\textheight]{example-image}{(d) 本章方法结果}
\end{minipage}\par\endgroup
\DefenseCaption{图3-2\quad 示例对比结果}
\DefenseTakeaway{本章方法的结果与参考值最接近}
\end{frame}

% layout: table
\begin{frame}{3.4 案例研究与结果分析}
\DefenseSubsection{3.4.2 示例指标对比}
\DefenseCaption{表3-1\quad 示例方法指标对比}
\begin{DefenseSource}
\centering
\adjustbox{max width=\textwidth,max totalheight=0.62\textheight}{%
\begin{tabular}{lccc}
  \toprule
  \multirow{2}{*}{方法} & \multicolumn{2}{c}{误差指标} & \multirow{2}{*}{\makecell{耗时\\(s)}}\\
  \cmidrule(lr){2-3}
  & RMSE & MAE & \\
  \midrule
  方法 A\cite{demo-ref} & 0.52 & 0.41 & 1.2\\
  方法 B & 0.47 & 0.36 & 1.5\\
  本章方法 & \textbf{0.39} & \textbf{0.30} & 1.4\\
  \bottomrule
\end{tabular}}\par
\end{DefenseSource}
\DefenseTakeaway{表中本章方法两项误差指标最低（合成数据）}
\end{frame}

% layout: paper-summary
\begin{frame}{3.5 本章小结}
\begin{itemize}
  \item 提出示例方法一，给出模型结构与训练目标。
  \item 在合成数据上验证了方法的有效性。
  \item 分析了方法的计算耗时与适用范围。
\end{itemize}
\DefensePaperBox{示例作者, 示例合作者. 示例方法一研究[J]. 示例期刊, 2026, 1(1): 1-10.}
\DefenseTakeaway{示例方法一回答了问题一}
\end{frame}

% layout: cards
\begin{frame}{总结与展望}
\DefenseSubsection{01\quad 主要创新点}
\DefenseCard{创新点 1}{针对示例数据的\DefenseHighlight{分布差异}问题，提出示例方法一，实现跨场景建模。（第 3 章）}
\DefenseCard{创新点 2}{针对标注样本\DefenseHighlight{稀缺}问题，构建示例方法二，降低标注需求。（第 4 章）}
\DefenseCard{创新点 3}{针对模型\DefenseHighlight{在线更新}问题，设计示例更新机制，保持部署精度。（第 5 章）}
\end{frame}

% layout: outlook
\begin{frame}{总结与展望}
\DefenseSubsection{02\quad 未来研究展望}
\DefenseBoxTitle{其一}\quad 将示例方法推广到更多数据类型，检验方法的适用边界。

\vspace{0.3cm}
\DefenseBoxTitle{其二}\quad 研究模型更新过程中的计算开销，降低在线部署成本。
\end{frame}

% layout: thanks
\begin{frame}[plain]
\DefenseThanksFrame
\end{frame}

\end{document}
